\begin{comment}
Este proyecto se enfoca en diseñar y fabricar un sistema de fabricación aditiva tipo
FDM con enfoque académico. 

El control numérico computarizado ha estado en constante desarrollo en la industria manufacturera. Actualmente se ha popularizado en equipos FDM como
lo son las impresoras 3D, que han tenido adopción en el uso domestico. 
La industria manufacturera se centra en la automatización total y fabricación flexible, y estos equipos FDM han fomentado el crecimiento de esta 
tecnología en entornos educativos.

El objetivo es proporcionar a estudiantes una herramienta educativa, practica y asequible para poner en practica conocimientos adquiridos en el 
transcurso de su estadía en la licenciatura. Estos equipos FDM son totalmente modificables por el usuario, lo cual convierte a este proyecto en uno con alta
oportunidad de mejora constante respecto a equipos comerciales, que son sistemas mas cerrados. Este trabajo consta de objetivos
específicos como el diseño y armado, las piezas serán impresas con otro equipo FDM siguiendo la filosofía Reprap [1].
Se usó software libre "Marlin", el cual es un firmware dedicado a equipos de impresión 3D. Se realizó
calibración de los módulos que conforman el equipo, ajuste final y verificación del comportamiento. Asi también, se diseñó 
el sistema de enfriamiento del extrusor y fue comparado con uno comercial, tanto en eficiencia de enfriamiento y calidad de impresión final.
Se obtuvieron resultados superiores en ciertos aspectos a equipos comerciales, tanto en calidad como velocidad de impresión.
En conclusión, se logró un diseño que tiene mejoras con un equipo comercial, que tiene ventajas como ser completamente modificable y ser un sistema mas abierto,
lo que permite la integración en un entorno académico.
\end{comment}

Este proyecto presenta el diseño y fabricación de una impresora 3D basada en tecnología de Modelado por Deposición Fundida (FDM) orientada a aplicaciones académicas y de prototipado. El desarrollo del sistema se fundamentó en una arquitectura de hardware abierto inspirada en la filosofía RepRap, permitiendo la fabricación de múltiples componentes mediante manufactura aditiva y facilitando la modificación mecánica, electrónica y de firmware.

El diseño mecánico fue desarrollado mediante modelado CAD paramétrico en SolidWorks, mientras que el control electrónico se implementó utilizando una plataforma basada en Arduino Mega 2560, electrónica RAMPS 1.4 y firmware Marlin. Como parte del proyecto, se rediseñó el sistema de extrusión, logrando una reducción del 24.2\% en su masa, y se desarrolló un sistema de enfriamiento de capa que permitió reducir la temperatura del disipador del hotend en un 42.25\%. Posteriormente, se realizaron procesos de ensamblaje, calibración y ajuste experimental de parámetros de impresión, incluyendo la sintonización del controlador térmico PID del hotend. 

El desempeño del sistema se evaluó mediante pruebas de precisión dimensional 
(obteniendo errores inferiores al 0.3\%), pruebas de torre de temperatura para 
determinar la ventana térmica óptima (205~$^\circ$C), y el diagnóstico 
experimental de absorción de humedad en el filamento. Finalmente, se concluye 
que el sistema representa una plataforma funcional, modular y de bajo costo 
(\$4189 MXN), con alto potencial como herramienta didáctica para la enseñanza 
de manufactura aditiva, control numérico y sistemas embebidos.

